\documentclass[pdf,bookmarks,colorlinks=true]{IEEEtran}
\usepackage{times}
\usepackage{amsmath}
\usepackage{hyperref}
\usepackage{url}
\usepackage{graphicx}
\usepackage{wrapfig}  

\title{\bf Generic Technical Paper Skeleton}

\author{Scot Anderson\\
Southern Adventist University\\
scot@southern.edu
}

\begin{document}

\maketitle

\begin{abstract}

The abstract should describe the basic message of the paper, including the problem, why your solution should be of interest, some notion that your solution is effective, and a teaser about how it has been evaluated. Cover all of this using between 75 and 150 words. The abstract is thus the hardest part to write. Sometimes I try to write it first, but the final version is usually composed of items drawn from the introduction, and then condensed, as the last step of writing the paper.

\end{abstract}

\section{Introduction}
\label{sec:Introduction}

\textbf{The problem we have solved:}

Concentrate on making {\em this} assertion and {\em only} this assertion in a succinct set of 1 to 3 paragraphs

A common mistake is to explain too much of the problem context first. Instead, state the problem essentially as a claim, and leave explanations supporting your claim to the next part, ``Why it is not already solved.''


\textbf{Why the problem is not already solved or other solutions are ineffective in one or more important ways}

Your new idea need not solve every problem but it should solve at least one that is not already solved

This is the place to provide a succinct description of the problem context giving enough information to support the claim that a problem exists, made in the preceding problem declaration.

\textbf{Why our solution is worth considering and why is it effective in some way that others are not}

A succinct statement of {\em why} the reader should care enough to read the rest of the paper.

This should include a statement about the characteristics of your solution to the problem which 1) make it a solution, and 2) make it superior to other solutions to the same problem.

\textbf{How the rest of the paper is structured}

The short statement below is often all you need, but you should change it when your paper has a different structure, or when more information is {\em required} to describe what a given section contains. If it isn't {\em required} then you don't want to say it here.

The rest of this paper first discusses related work in \ref{sec:RelatedWork}, and then describes our implementation in \ref{sec:Implementation}. \ref{sec:Evaluation} describes how we evaluated our system and presents the results. \ref{sec:Conclusion} presents our conclusions and describes future work.

\section{Related Work}\label{sec:RelatedWork}

\textbf{Other efforts that exist to solve this problem and why are they less effective than our method} Resist the urge to point out only flaws in other work. Do your best to point out both the strengths and weaknesses to provide as well rounded a view of how your idea relates to other work as possible 

In a social and political sense, it is {\em very smart} as well as ethical to say good things, which are true, about other people's work. A major motivation for this is that editors and program committee members have to get a set of reviews for your paper. The easiest way for them to decide who should review it is to look at the set of references to {\em related work} (e.g., \cite{Anderson2007,Anderson2009}) to find people who are likely to be competent to review your paper.  The people whose work you talk about are thus likely to be reading what you say about {\em their} work while deciding what to say about {\em your} work.

Clear enough? Speak the truth, say what you have to say, but be generous to the efforts of others.

\textbf{Other efforts that exist to solve related problems that are relevant, how are they relevant, and why are they less effective than our solution for this problem}
Many times no one has solved your exact problem before, but others have solved closely related problems or problems with aspects that are strongly analogous to aspects of your problem.

\section{Implementation}\label{sec:Implementation}

\textbf{What we (will do $|$ did): {\em Our Solution}} Another way to look at this section is as a paper, within a paper, describing your implementation. That viewpoint makes this the introduction to the subordinate paper, which should describe the overall structure of your implementation and how it is designed to address the problem effectively.

Then, describe the structure of the rest of this section, and what each subsection describes.


\textbf{How our solution (will $|$ does) work} This is the body of the subordinate paper describing your solution. It may be divided into several subsections as required by the nature of your implementation.

The level of detail about how the solution works is determined by what is appropriate to the type of paper (conference, journal, technical report)

This section can be fairly short for conference papers, fairly long for journal papers, or {\em quite} long in technical reports. It all depends on the purpose of the paper and the target audience

Proposals are necessarily a good deal more vague in this section since you have to convince someone you know enough to have a good chance of building a solution, but that you have not {\em already} done so.

\section{Evaluation}\label{sec:Evaluation}

\textbf{How we tested our solution}
\begin{itemize}
\item   Performance metrics
\item   Performance parameters
\item   Experimental design
\end{itemize}

\textbf{How our solution performed, how its performance compared to that of other solutions mentioned in related work, and how these results show that our solution is effective}:

\begin{itemize}
\item   Presentation and Interpretation
\item   Why, how, and to what degree our solution is better
\item   Why the reader should be impressed with our solution
\item   Comments
\end{itemize}

\textbf{Context and limitations of our solution as required for summation:} State what the results {\em do} and {\em do not} say.

\section{Conclusions and Future Work}\label{sec:Conclusion}

\textbf{The problem we have solved}: The most succinct statement of the problem in the paper. Ideally one sentence. More realistically two or three. Remember that you simply state it without argument. If you have written a good paper you are simply reminding the reader of what they now believe and of how much they agree with you.

\textbf{Our solution to the problem}: Again, a succinct statement of the presented solution. Sometimes it works well to leave it at that and not even describe your solution here. If you do, then again state your solution in one or two sentences taking the rhetorical stance that this is all obvious. If you have a good solution and have written an effective paper, then the reader already agrees with you.

\textbf{Why our solution is worthwhile in some significant way}: Again, a succinct restatement in just a few sentences of why your solution is worthwhile assuming the reader already agrees with you.

\textbf{Why the reader should be impressed and/or pleased to have read the paper}: A few sentences about why your solution is valuable, and thus why the reader should be glad to have read the paper and why they should be glad you did this work.

\textbf{What we will (or could) do next}

\begin{itemize}
\item   Improve our solution
\item   Apply our solution to harder or more realistic versions of this problem
\item   Apply our solution or a related solution to a related problem
\end{itemize}

\bibliographystyle{IEEETrans}
\bibliography{./my}

\appendix
\section{Using Tables and Graphics}

A picture is worth a thousand words and a Table is worth at least one hundred. Consequently, you should use images and tables as often as possible to support your text. The following tables and images provide examples.

\subsection{Images}

Images should always be in a vector scalar format when possible, and it is nearly always possible. Almost all scalar graphics programs can export the image to a portable document format (PDF). Figure~\ref{fig:crow} shows how to  easily include PDF files. This PDF started life as a scalar vector graphic (SVG) file create in InkScape. 

\begin{figure}[h]
    \centering
    \includegraphics[trim={0.1cm 1cm 0.1cm 1cm},clip,width=\columnwidth]{figures/Crow.pdf}
    \caption{Original artwork by Scot Anderson \copyright 2018.}
    \label{fig:crow}
\end{figure}

The shape of your image or table may affect how you choose to display it. For example, a nice portrait image, such as Figure~\ref{fig:crow}, fits nicely in a column. Tables rarely fit in columns, hence we may need to display them across columns. To accomplish this with tables or figures add a star (*) to the word figure or table. 

Table~\ref{tbl:frequencies} shows the ranges of radio frequencies.

\begin{table}[htb]
    \centering
    \caption{Radio Spectrum.}
    \label{tbl:frequencies}
    \begin{tabular}{|l|l|l|l|}
        \hline
        Name & Abbreviation & Frequency & Wavelength \\
        \hline \hline
        Very low frequency & VLF & 3-30 kHz & 100 - 10 km \\ \hline
        Low frequency & LF & 30-300 kHz & 10 - 1 km \\ \hline
        %Medium frequency & MF & 300-3,000 kHz & 1,000 - 100 m \\ \hline
        %High frequency & HF & 3-30 MHz  & 100 - 10 m \\ \hline
        %Very high frequency & VHF & 30-300 MHz & 10 - 1 m \\ \hline        
    \end{tabular}
\end{table}

\end{document}
